
\documentclass[11pt,a4paper]{article}

% ==================================================
% 1. PACKAGES
% ==================================================

\usepackage[a4paper, margin=1in]{geometry}
\usepackage{amsmath, amssymb, amsthm, mathtools}
\usepackage{mathrsfs}
\usepackage{bm}
\usepackage{enumitem}
\usepackage{xcolor}
\usepackage[most]{tcolorbox}
\usepackage{hyperref}
\usepackage{cleveref}
\usepackage{mathtools}

\hypersetup{
    colorlinks=true,
    linkcolor=blue,
    citecolor=blue,
    urlcolor=blue
}

\setlength{\parskip}{0.35em}
\setlength{\parindent}{0pt}

% ==================================================
% 2. MATHEMATICAL ENVIRONMENTS
% ==================================================

\theoremstyle{definition}
\newtheorem{definition}{Definition}[section]
\newtheorem{example}[definition]{Example}
\newtheorem{exercise}[definition]{Exercise}
\newtheorem{remark}[definition]{Remark}
\newtheorem{counterexample}[definition]{Counterexample}
\newtheorem{note}[definition]{Note}

\theoremstyle{plain}
\newtheorem{theorem}[definition]{Theorem}
\newtheorem{lemma}[definition]{Lemma}
\newtheorem{proposition}[definition]{Proposition}
\newtheorem{corollary}[definition]{Corollary}

% ==================================================
% 3. PERSONAL NOTE BOX
% ==================================================

\newtcolorbox{mynote}[1][Personal Note]{
    enhanced,
    breakable,
    colback=yellow!5,
    colframe=yellow!50!black,
    title={#1},
    fonttitle=\bfseries,
    boxrule=0.6pt,
    arc=2mm,
    left=2mm,
    right=2mm,
    top=1mm,
    bottom=1mm
}

% ==================================================
% 4. CUSTOM MATHEMATICAL COMMANDS
% ==================================================

\newcommand{\R}{\mathbb{R}}
\newcommand{\N}{\mathbb{N}}
\newcommand{\Q}{\mathbb{Q}}
\newcommand{\Z}{\mathbb{Z}}
\newcommand{\E}{\mathbb{E}}
\newcommand{\A}{\mathcal{A}}
\newcommand{\Prob}{\mathbb{P}}
\newcommand{\eps}{\varepsilon}

\newcommand{\norm}[1]{\left\lVert #1 \right\rVert}
\newcommand{\abs}[1]{\left| #1 \right|}
\newcommand{\inner}[2]{\left\langle #1,#2 \right\rangle}
\newcommand{\set}[1]{\left\{#1\right\}}

% ==================================================
% 5. DOCUMENT INFORMATION
% ==================================================

\title{\textbf{Notations}}

% ==================================================
% 6. DOCUMENT
% ==================================================

\begin{document}

\maketitle
\tableofcontents
\newpage

\section{Useful facts}

\begin{theorem}[Sum of arbitary set of numbers]
We have the following:
\[
\sum_{\alpha\in A}x_{\alpha} = \sup\set{\sum_{F\subset A,\alpha \in F}x_{\alpha} \text{ and } F \text{ finite}}.
\]
\end{theorem}


\begin{mynote}[Axiom of choice]
If we have infinite non-empty sets, there will always exist a way that we could take an element of each set.
\end{mynote}


\section{Exercises}

All exercises are from \textit{An Introduction to Measure Theory} by Terence Tao.

\begin{exercise}
\textbf{\textit{See Exercise 0.0.1 p.xiii.}}

\textbf{Proof idea}. The exercise wants us to proof that, when a sum of infinite (may uncountable) of non-negative numbers is finite, then that sequence has countable positive numbers.  We need to prove $U = \set{x_{\alpha} > 0: \alpha \in A}$ is countable. Of course if we try to count it or prove it countable directly, it gonna be hard. So we gonna divide $U$ to smaller sets, then prove those smaller sets are countable, then $U$ is countable.

\begin{proof}
Let $U_n = \set{\alpha\in A: x_{\alpha}\ge \frac{1}{n}}$.

Hence, we have $U_n \uparrow U$. (1)

Let the sum: $S = \sum_{\alpha \in A}x_\alpha = M < \infty$.

Here, we assume each $U_n$ is infinite for all $n$.

Then, we take a finite number in $U_n$, to try to calculate $S$, but $U_n$ is infinite, then we could take arbitrary large number in that, suppose we take $k$ numbers, then:
\[
\sum_{F\subset U_n, x_\alpha \in F}x_\alpha \ge \frac{1}{n}k.
\]

But we can take arbitrary large $k$, then we choose $k = (nM + 3)$, then we have:

\[
\sum_{F\subset U_n, x_\alpha \in F}x_\alpha \ge \frac{1}{n}(nM+3) = M + \frac{3}{n}>M.
\]

but:

\[
\sum_{F\subset U_n, x_\alpha \in F}x_\alpha < M.
\]
(by using Theorem 1.1).

Then, it is a contradiction, $U_n$ must be finite. (2)

From (1) and (2), $U$ must be finite, hence it is countable.
\end{proof}
\end{exercise}

\begin{exercise}

\textbf{\textit{This is the proof for Theorem 0.0.2 p.xiii.}}
\begin{proof}
\textbf{Prove $\sum_{(n,m)\in \N^2}x_{n,m}\le \sum_{n=1}^{\infty}\sum_{m=1}^{\infty}x_{n,m}$}:

We have $\sum_{(n,m)\in \N^2}x_{n,m} = \sup{\sum_{F\subset \N^2, F \text{ finite}, (n,m)\in F}x_{n,m}}$ by Theorem 1.1.

Take a random finite $F$, then we can always find $M$ and $N$ such that $M= \max\set{m:(n,m)\in F}$ and $N= \max\set{n:(n,m)\in F}$. After that, construct a rectangle: $\set{1,2,3,...,N} \times \set{1,2,3,...,M}$ We have:
\[
\sum_{(n,m)\in F}x_{n,m} \le \sum_{n=1}^{N}\sum_{m=1}^{M}x_{n,m} \le \sum_{n=1}^{\infty}\sum_{m=1}^{\infty}x_{n,m}.
\]
for all finite $F$. Then:

\[
\sup\set{\sum_{(n,m)\in F}x_{n,m}} = \sum_{(n,m)\in \N^2}x_{n,m}  \le \sum_{n=1}^{\infty}\sum_{m=1}^{\infty}x_{n,m}.
\]


\textbf{Prove $\sum_{(n,m)\in \N^2}x_{n,m}\ge \sum_{n=1}^{\infty}\sum_{m=1}^{\infty}x_{n,m}$}:

With any finite $M = 1,2,3...$ and $N = 1,2,3,...$ we have:

\[
\sum_{n=1}^{N}\sum_{m=1}^{M}x_{n,m} \le \sum_{(n,m)\in \N^2}x_{n,m} 
\]
(by using Theorem 1.1), then if we take $\lim$ on the left hand side, we will obtain the inequality.

\textbf{Prove $\sum_{n=1}^{\infty}\sum_{m=1}^{\infty}x_{n,m} = \sum_{m=1}^{\infty}\sum_{n=1}^{\infty}x_{n,m}$}

Very same concept, for finite, we can always interchange the index, but in the infinite case, you need to take $\lim$ on both side. Note that, the order of taking $\lim$ matters.
\end{proof}


\end{exercise}

\end{document}
